% Volume III: Attention and Publicity
% Part V. Publicity Without Comprehension
% Chapter 28: Collective Investigation
% Spine: proof-spines/ch028-spine.md

\chapter{Collective Investigation}
\label{ch:ch028}

Large crowds can parallelize discovery, but they also generate shared priors, copied mistakes, harassment, false identification, and narrative cascades. The chapter models collective investigation as a distributed search process with non-independent errors. Its effective number of independent investigators is not the head count \(n\), but approximately
\[
n_{\mathrm{eff}} = \frac{n}{1+(n-1)\rho},
\]
\ToyModel\ where \(\rho\) is average error correlation. As \(\rho\) rises, crowds become less like many checks and more like one amplified misreading.

The argument therefore favors organized collective inquiry rather than romantic crowd omniscience. Reliable public knowledge requires institutions and norms that lower \(\rho\): role differentiation, adversarial checking, preserved records, and explicit thresholds for moving from clue to claim. Publicity becomes epistemically useful only when scattered observers can be coordinated without being collapsed into a single cascade. Those same thresholds become unmeetable burdens when they are used to discount dispersed but convergent clues until only already-authorized investigators can move from suspicion to claim.
